\documentclass[11pt]{article}
\usepackage[margin=1in]{geometry}
\usepackage{amsmath,amssymb,amsthm,mathtools}
\usepackage{graphicx}
\usepackage[hidelinks]{hyperref}
\usepackage{microtype}

\newtheorem{theorem}{Theorem}
\newtheorem{lemma}[theorem]{Lemma}
\newtheorem{corollary}[theorem]{Corollary}
\newcommand{\B}{\mathcal B}
\newcommand{\N}{\mathcal N}
\newcommand{\R}{\mathcal R}
\newcommand{\UT}{\mathsf{UT}}

\title{Invertible Completions of Upper-Triangular Operator Matrices\\
at Every Finite Matrix Size}
\author{Candidate full solution to the higher-dimensional question in arXiv:2202.04960}
\date{17 August 2026}

\begin{document}
\maketitle

\begin{center}
\fbox{\parbox{0.91\linewidth}{\textbf{Status.} Candidate full solution,
likely valid pending expert review.  The criterion is necessary and sufficient
for every finite matrix size and arbitrary Banach spaces, assuming only that
the interior diagonal operators are inner regular.}}
\end{center}

\section{Source question}

Let $X_1,\ldots,X_n$ be Banach spaces, let $D_i\in\B(X_i)$ be fixed,
and let
\[
 T_n^d(A)=
 \begin{pmatrix}
 D_1&A_{12}&\cdots&A_{1n}\\
 0&D_2&\cdots&A_{2n}\\
 \vdots&\ddots&\ddots&\vdots\\
 0&\cdots&0&D_n
 \end{pmatrix},
 \qquad A_{ij}\in\B(X_j,X_i),\quad i<j.
\]
Sarajlija and Djordjevi\'c give a $3\times3$ completion theorem on arbitrary
Banach spaces.  In Section 3 of \cite{SD}, PDF page 10, they ask whether the
$2\times2$ and $3\times3$ results admit a natural generalization to all
$n>3$, anticipating multiple Fredholm-type assumptions and several
left-invertible defect maps.  The complete passage is reproduced below.

\begin{center}
\includegraphics[width=0.96\linewidth]{figures/open_problem_crop.png}
\end{center}

The theorem below answers this question by isolating the exact recursive
defect condition.  It also sharpens the source's Theorem 2.1: its condition
1 is not only sufficient but necessary for an invertible $3\times3$
completion.

\section{Proof intuition}

An inner regular operator $D_i$ splits into an isomorphism from a complement
of $\N(D_i)$ onto $\R(D_i)$, together with a zero map from its kernel to its
cokernel.  After all invertible pieces are collected, ordinary Schur
elimination reduces the entire completion problem to a single operator
between defect sums:
\[
 \N(D_2)\oplus\cdots\oplus\N(D_n)
 \longrightarrow
 X_1/\R(D_1)\oplus\cdots\oplus X_{n-1}/\R(D_{n-1}).
\]
The original zero pattern says exactly that this defect operator is upper
triangular.

The remaining question is purely about Banach spaces: when does an invertible
upper-triangular operator exist between two finite direct sums?  Its first
diagonal block must be left invertible and its last diagonal block must be
right invertible.  Split off those two matched pieces.  Gaussian row and
column operations leave the same problem with one fewer row and column.
Iterating this operation gives a finite, exact recursive criterion.  At
$n=3$ the last residual condition is precisely the quotient isomorphism
appearing in the source paper.

\section{The recursive defect relation}

All direct sums below are finite topological direct sums; their particular
standard norm is irrelevant up to isomorphism.  For two lists of Banach
spaces $(E_1,\ldots,E_m)$ and $(F_1,\ldots,F_m)$, write
\[
 \UT(E_1,\ldots,E_m;F_1,\ldots,F_m)
\]
if there is an invertible upper-triangular block operator
\[
 U:\bigoplus_{j=1}^m E_j\longrightarrow\bigoplus_{i=1}^m F_i,
 \qquad U_{ij}=0\quad(i>j).
\]

\begin{lemma}[recursive triangular-isomorphism criterion]\label{lem:rec}
For $m=1$, $\UT(E_1;F_1)$ holds exactly when $E_1\cong F_1$.
For $m\geq2$, $\UT(E_1,\ldots,E_m;F_1,\ldots,F_m)$ holds exactly when there
exist left-invertible operators
\[
 L:E_1\to F_1,\qquad J:F_m\to E_m
\]
such that
\begin{equation}\label{eq:recursion}
 \UT\left(E_2,\ldots,E_{m-1},E_m/J(F_m);
 F_1/L(E_1),F_2,\ldots,F_{m-1}\right)
\end{equation}
holds.  The ranges of $L$ and $J$ are automatically closed and complemented,
so the displayed quotients are Banach spaces.
\end{lemma}

\begin{proof}
Suppose first that $U=(U_{ij})$ is invertible and upper triangular.  Its first
column is $(U_{11},0,\ldots,0)^t$.  Since $U$ is bounded below,
$L:=U_{11}:E_1\to F_1$ is bounded below.  More precisely,
\[
 F_1\xrightarrow{\iota_1}\bigoplus_iF_i
 \xrightarrow{U^{-1}}\bigoplus_jE_j
 \xrightarrow{\pi_1}E_1
\]
is a bounded left inverse of $L$.  Thus $L$ is left invertible.  Surjectivity
of the last row of $U$ shows that $R:=U_{mm}:E_m\to F_m$ is onto, and
\[
 F_m\xrightarrow{\iota_m}\bigoplus_iF_i
 \xrightarrow{U^{-1}}\bigoplus_jE_j
 \xrightarrow{\pi_m}E_m
\]
is a bounded right inverse $J$ of $R$.  In particular, $J:F_m\to E_m$ is
left invertible and
\[
 F_1=L(E_1)\oplus F_1',\qquad
 E_m=E_m'\oplus J(F_m)
\]
for closed complements $F_1'$ and $E_m'$.

Use a triangular domain operation to subtract from the $E_1$ coordinate the
component that maps the remaining columns into $L(E_1)$.  Use a triangular
range operation to subtract from the earlier rows the components produced by
$J(F_m)$, whose last-row image is all of $F_m$.  Both operations are
invertible.  Relative to the displayed splittings they transform $U$ into
\[
 \begin{pmatrix}
 L_0&0&0\\
 0&U'&0\\
 0&0&R_0
 \end{pmatrix},
\]
where $L_0:E_1\to L(E_1)$ and $R_0:J(F_m)\to F_m$ are isomorphisms, and
\[
 U':E_2\oplus\cdots\oplus E_{m-1}\oplus E_m'
 \longrightarrow F_1'\oplus F_2\oplus\cdots\oplus F_{m-1}
\]
is still upper triangular.  Since $U$ is invertible, $U'$ is invertible.
Identifying the complements with the corresponding quotients proves
\eqref{eq:recursion}.

Conversely, choose complements as above and let $U'$ realize
\eqref{eq:recursion}.  Define $U$ to be $L_0$ on $E_1$, $R_0=J^{-1}$ on
$J(F_m)$, and $U'$ on the remaining summands, with every other block zero.
This operator is upper triangular in the original lists and is the direct sum
of three isomorphisms.  Hence it is invertible.
\end{proof}

\section{Defect reduction}

Recall that $D\in\B(X)$ is inner regular precisely when its kernel and range
are closed and complemented.  In that case choose
\[
 X=\N(D)\oplus M,\qquad X=\R(D)\oplus K.
\]
Then $D|_M:M\to\R(D)$ is an isomorphism and
$K\cong X/\R(D)$.

\begin{lemma}[Schur reduction to the defects]\label{lem:defect}
Assume $D_1,\ldots,D_n$ are inner regular.  There is a choice of the
off-diagonal blocks $A_{ij}$ for which $T_n^d(A)$ is invertible if and only if
there is an invertible block operator
\begin{equation}\label{eq:strict}
 S:\bigoplus_{j=1}^n\N(D_j)\longrightarrow
   \bigoplus_{i=1}^n X_i/\R(D_i),
 \qquad S_{ij}=0\quad(i\geq j).
\end{equation}
Moreover every strict upper-triangular $S$ in \eqref{eq:strict} can be
realized by a choice of the $A_{ij}$.
\end{lemma}

\begin{proof}
Choose complements $M_i,K_i$ as above and identify $K_i$ with
$X_i/\R(D_i)$.  Reorder the domain as
$M\oplus N=(\bigoplus_iM_i)\oplus(\bigoplus_i\N(D_i))$ and the range as
$R\oplus K=(\bigoplus_i\R(D_i))\oplus(\bigoplus_iK_i)$.  In these
coordinates any completion has the form
\[
 \widetilde T=\begin{pmatrix}P&Q\\ V&W\end{pmatrix}:M\oplus N\to R\oplus K.
\]
The block $P:M\to R$ is upper triangular and its diagonal blocks are the
isomorphisms $D_i|_{M_i}$.  Hence $P$ is invertible.  Schur elimination gives
\[
 \begin{pmatrix}I&0\\-VP^{-1}&I\end{pmatrix}
 \widetilde T
 \begin{pmatrix}I&-P^{-1}Q\\0&I\end{pmatrix}
 =\begin{pmatrix}P&0\\0&W-VP^{-1}Q\end{pmatrix}.
\]
Thus $\widetilde T$ is invertible exactly when
$S:=W-VP^{-1}Q:N\to K$ is invertible.

The original block pattern permits a map from the $j$th input to the $i$th
output only when $i\leq j$, while $D_j$ vanishes on $\N(D_j)$.  Every path
contributing to $VP^{-1}Q$ has strictly increasing source-to-target index.
Consequently $S_{ij}=0$ whenever $i\geq j$.

Conversely, given such an $S$, define $A_{ij}$ to equal the composition
\[
 X_j\longrightarrow\N(D_j)\xrightarrow{S_{ij}}K_i\hookrightarrow X_i
\]
for $i<j$, using the bounded projection associated with the chosen domain
splitting, and set all other components of $A_{ij}$ to zero.  Then
$Q=V=0$, $P=\bigoplus_iD_i|_{M_i}$, and the Schur complement is exactly
$S$.
\end{proof}

\section{Full all-dimensional characterization}

\begin{theorem}[all-$n$ invertible completion criterion]\label{thm:main}
Let $n\geq2$, let $X_1,\ldots,X_n$ be arbitrary Banach spaces, and let
$D_i\in\B(X_i)$.  Assume only that the interior operators
$D_2,\ldots,D_{n-1}$ are inner regular.  Put
\[
 E_r=\N(D_{r+1}),\qquad F_r=X_r/\R(D_r),qquad 1\leq r\leq n-1.
\]
Then the following are equivalent:
\begin{enumerate}
\item There are blocks $A_{ij}\in\B(X_j,X_i)$, $i<j$, for which
$T_n^d(A)$ is invertible.
\item $D_1$ is left invertible, $D_n$ is right invertible, and
\[
 \boxed{\ \UT(E_1,\ldots,E_{n-1};F_1,\ldots,F_{n-1})\ }
\]
holds, equivalently the finite recursion in Lemma~\ref{lem:rec} terminates in
an isomorphism.
\end{enumerate}
\end{theorem}

\begin{proof}
Suppose $T_n^d(A)$ is invertible.  Its restriction to the first coordinate
shows that $D_1$ is bounded below and has complemented range: a bounded left
inverse is obtained from the first coordinate of $T_n^d(A)^{-1}$.  Thus
$D_1$ is left invertible.  Applying $T_n^d(A)^{-1}$ to vectors supported in
the last output coordinate gives a bounded right inverse for $D_n$.  Hence
$D_n$ is right invertible.  In particular both endpoint operators are inner
regular, so Lemma~\ref{lem:defect} applies.

In \eqref{eq:strict}, the first domain summand $\N(D_1)$ maps to zero and the
last range summand $X_n/\R(D_n)$ receives only zero.  Invertibility therefore
forces both spaces to vanish, as already follows from endpoint one-sided
invertibility.  Deleting them turns \eqref{eq:strict} into an ordinary
upper-triangular isomorphism
\[
 \bigoplus_{r=1}^{n-1}E_r\longrightarrow\bigoplus_{r=1}^{n-1}F_r.
\]
This is exactly the displayed $\UT$ condition.

Conversely, endpoint one-sided invertibility and interior regularity make all
$D_i$ inner regular.  A $\UT$ witness gives the strict defect isomorphism
in \eqref{eq:strict} after adjoining the zero endpoint defect spaces.
Lemma~\ref{lem:defect} realizes it by off-diagonal blocks and proves the
resulting completion invertible.
\end{proof}

\begin{corollary}[$3\times3$ strengthening]\label{cor:three}
Under the hypotheses of Theorem 2.1 of \cite{SD}, an invertible completion
exists if and only if $D_1$ is left invertible, $D_3$ is right invertible,
and there are left-invertible maps
\[
 J_1:\N(D_2)\to X_1/\R(D_1),\qquad
 J_2:X_2/\R(D_2)\to\N(D_3)
\]
such that
\[
 \frac{X_1/\R(D_1)}{J_1\N(D_2)}
 \cong
 \frac{\N(D_3)}{J_2(X_2/\R(D_2))}.
\]
Thus condition 1 of that theorem is equivalent to its completion condition
2, not merely sufficient for it.
\end{corollary}

\begin{proof}
Apply Lemma~\ref{lem:rec} with $m=2$ to Theorem~\ref{thm:main}.
\end{proof}

For $n=2$, Theorem~\ref{thm:main} reduces to the classical criterion
$X_1/\R(D_1)\cong\N(D_2)$ of Han, Lee, and Lee \cite{HLL}.  For finite
dimensional defect spaces, Lemma~\ref{lem:rec} becomes the usual sequence of
prefix dimension inequalities together with equality of total dimensions.

\section{Verification, scope, and novelty}

The accompanying script constructs a nontrivial upper-triangular defect
isomorphism with range-block dimensions $(3,1,1)$ and domain-block dimensions
$(1,1,3)$, embeds it as the Schur complement of an invertible reduced block
matrix, and checks a dimension pattern with an obstructed first recursive
embedding.  It also verifies the Schur determinant identity and strict zero
pattern for random off-diagonal blocks.  These checks are not part of the
proof.  (A genuinely nontrivial one-sided defect for a square diagonal block
is inherently infinite-dimensional, so ordinary finite square matrices cannot
model the whole completion theorem.)

The result is invariant under all choices of kernel/range complements: the
statement uses canonical quotient spaces, and different complements are
isomorphic to them.  No separability, Hilbert structure, or Fredholm
finite-dimensionality is used.  The matrix size must be finite; infinite
nests would require convergence and uniform-boundedness hypotheses absent
from the source question.

A bounded novelty check on 17 August 2026 covered the run's four lightweight
indexes, the local full sources for arXiv:2108.12153, 2110.07387, 2112.02350,
and 2202.04960, and exact-title/core-keyword web searches through that date.
Theorem 2.13 of \cite{S21}, cited by the source, gives a sufficient condition
requiring adjacent defect spaces to be isomorphic and a weaker Hilbert-space
necessity statement.  The other located papers concern spectral stability or
Fredholm/Weyl completion, rather than the exact ordinary-invertibility
criterion above.  No recursive necessary-and-sufficient Banach-space result,
nor the $3\times3$ necessity upgrade in Corollary~\ref{cor:three}, was
located.  Novelty confidence is moderate pending a specialist search,
especially in the block-operator monograph literature.

\section{Human review recommendation}

First check that the Schur complement in Lemma~\ref{lem:defect} retains the
strict index pattern; this is the bridge from operator completion to defect
matching.  Then check the two Gaussian eliminations in Lemma~\ref{lem:rec},
including preservation of upper triangularity.  Finally compare
Corollary~\ref{cor:three} line by line with condition 1(c) of the source's
Theorem 2.1.  Those checks settle both the all-$n$ result and the claimed
$3\times3$ strengthening.

\begin{thebibliography}{9}
\bibitem{SD}
N. Sarajlija and D. S. Djordjevi\'c,
\emph{Completion problem of upper triangular $3\times3$ operator matrices
on arbitrary Banach spaces}, arXiv:2202.04960v4 (2025), published in
Bull. Iranian Math. Soc. 51 (2025), Paper 85.

\bibitem{S21}
N. Sarajlija,
\emph{Perturbing the spectrum of operator $T_n^d(A)$},
arXiv:2108.12153; RACSAM 117 (2023), Paper 10, especially Theorem 2.13.

\bibitem{HLL}
J. K. Han, H. Y. Lee, and W. Y. Lee,
\emph{Invertible completions of $2\times2$ upper triangular operator
matrices}, Proc. Amer. Math. Soc. 128 (2000), 119--123.
\end{thebibliography}

\end{document}
